\section{Análisis de casos de la ventana de producto: del asombro a la herramienta predeterminada}

La sección anterior presentó la metáfora de la minería; esta sección la desglosa en el ciclo de vida de un producto para responder a la pregunta «¿por qué el Magic Moment deja de bastar tan pronto?». La trayectoria habitual de un producto de IA es la siguiente: surge una nueva capacidad de un modelo, un emprendedor la empaqueta en una primera experiencia que asombra, los usuarios la difunden con rapidez, la empresa del modelo y los competidores la replican y, al final, esa capacidad se convierte en una función predeterminada. Solo si, después de la difusión, el producto sigue acumulando flujos de trabajo, datos y marca, evita ser absorbido por la suite integral de la plataforma.

\begin{knowledgebox}{Ciclo de vida del Magic Moment}
\begin{tabularx}{\textwidth}{p{0.18\textwidth}p{0.34\textwidth}X}
\toprule
Etapa & Qué ocurre & Acción clave \\
\midrule
Descubrimiento de la capacidad & La nueva capacidad del modelo acaba de estar disponible & Encontrar pronto la tarea que mejor demuestre su valor. \\
Empaquetado de la experiencia & Se convierte en un Magic Moment que el usuario puede percibir & Reducir la barrera de uso y reforzar lo que impulsa la difusión. \\
Difusión explosiva & Se propaga por redes sociales y de boca en boca & Ocupar el lugar predeterminado en la mente del usuario y construir una marca temprana. \\
Reacción de la competencia & La empresa del modelo o productos similares la copian & Completar flujos de trabajo, datos y colaboración en equipo. \\
Conversión en función predeterminada & La capacidad pasa a ser una función básica de la plataforma & Migrar a escenarios más profundos o convertirse en un componente de la plataforma. \\
\bottomrule
\end{tabularx}
\end{knowledgebox}

\begin{importantbox}{La línea divisoria entre el asombro y la herramienta predeterminada}
El asombro nace de «ver por primera vez que puede hacerlo»; la herramienta predeterminada nace de «no poder prescindir de ella ningún día». Lo difícil para un producto de IA es completar esta transición antes de que se cierre la ventana del asombro.
\end{importantbox}

\begin{warningbox}{Efectos secundarios del Magic Moment}
El Magic Moment puede hacer creer al equipo que el producto ya está consolidado, pero el asombro del usuario no equivale a retención a largo plazo. El verdadero peligro es que el equipo siga optimizando en función del efecto de la demostración, en lugar de hacerlo en función de las tareas de alta frecuencia, los procesos de colaboración y las razones para pagar.
\end{warningbox}

\subsection{Resumen de la sección}

El ciclo de vida de los productos de IA es cada vez más corto, y los emprendedores deben llevar el producto de una experiencia que asombra a una herramienta predeterminada. De lo contrario, el Magic Moment no será más que una función dentro de la suite integral de la empresa del modelo.
